% TOP_PROBLEM: {"id":30007649,"problem_number":"LOCAL-30007649","title":"Partner sorting and inequality","category_id":21,"category":{"id":21,"name":"economics","display_name":"Economics","description":"Open economic theory statements, published research agendas, and proposed scoped specifications; question provenance and historical priority are qualified per record.","slug":"economics","order_index":21,"created_at":"2026-10-08T00:00:00Z"},"status":"provenance_pending","external_url":null,"legacy_ids":[],"published":false,"statement":"In the explicitly specified setting: How much do changing marriage-market opportunities and assortative mating causally amplify household inequality and intergenerational persistence? The mathematical statement above fixes the target, assumptions and scope of this catalogue formulation.","economics":{"original_id":138,"field":"Inequality and economic mobility","kind":"empirical","formalization_basis":"proposed_specification","source_status":"not_verified","priority_status":"not_established","priority_note":"No qualifying problem statement was directly verified, so historical priority cannot be assigned.","earliest_verified_reference":null,"current_reference":{"authors":["Benjamin Goldman","Jamie Gracie","Sonya R. Porter"],"title":"Who Marries Whom? The Role of Segregation by Race and Class","year":2026,"url":"https://www.nber.org/papers/w35140","doi":"10.3386/w35140","locator":"Abstract, paragraphs on childhood exposure and the spatial marriage-market model","evidence_summary":"The paper identifies exposure effects on marriage and quantifies segregation; the full downstream inequality and descendant-mobility problem is not explicitly posed as open."},"formalization_note":"This is a proposed scoped research specification. No primary passage explicitly stating this entire remaining question as open was verified. The supporting literature may answer important subquestions; this entry must not be advertised as a verified formal open problem."},"statusQualification":"Provenance pending: an explicit published statement of this exact open question was not verified. This is a catalogue-authored candidate specification, not a certified open conjecture. This is a proposed scoped research specification. No primary passage explicitly stating this entire remaining question as open was verified. The supporting literature may answer important subquestions; this entry must not be advertised as a verified formal open problem.","statusReviewedAt":"2026-10-08"}
% FIRST_POSTED: 2026-10-08
% ENTRY_KIND: statement_only
% !TeX program = lualatex
\documentclass[11pt]{article}
\usepackage{fontspec}
\usepackage[a4paper,margin=25mm]{geometry}
\usepackage{amsmath,amssymb,mathtools}
\usepackage{xurl}
\usepackage[hidelinks]{hyperref}
\newcommand{\sref}[2]{\hyperlink{source:#1:#2}{[S#2]}}
\setlength{\emergencystretch}{3em}
\begin{document}
% problemId:30007649
\section*{Partner sorting and inequality}\label{30007649}
\subsection{Definitions and mathematical statement}
Fix a cohort of potential partners and a feasible intervention \(z\) changing opportunities to meet across socioeconomic groups. Let \(M(z)\) be the equilibrium matching of partners, allowing some individuals to remain single. Person \(i\) has potential labor earnings \(y_i(M,z)\), reflecting endogenous labor supply as well as partner formation. For household \(h\), define equivalent disposable income \(Y_h(z)\) using a declared sharing and household-size rule. Let \(I(F_z)\) be a specified inequality functional of the household-income distribution \(F_z\), and let \(R_i(z)\) be the adult income rank of a subsequently born child under a stated cohort-selection rule. Estimate \(I(F_z)-I(F_0)\) and child mobility effects. Separate meeting opportunities from partner preferences by varying exposure without assuming preferences remain unchanged after exposure. Distinguish a mechanical reassignment exercise holding earnings fixed from a causal counterfactual that changes work, fertility, and investment. Identification must address selection into institutions, migration, marriage, and parenthood. The proposed problem is to quantify these downstream distributional consequences; evidence that segregation affects who marries whom does not automatically identify the full inequality and next-generation response.

\par\textbf{Formulation and scope.} This is a proposed scoped research specification. No primary passage explicitly stating this entire remaining question as open was verified. The supporting literature may answer important subquestions; this entry must not be advertised as a verified formal open problem.

\par\textbf{Open-question provenance.} An explicit published open-question statement was not verified. Sources below are supporting literature and must not be treated as a first listing.

\par\textbf{Historical priority.} No qualifying problem statement was directly verified, so historical priority cannot be assigned.

\subsection{Short English statement}
In the explicitly specified setting: How much do changing marriage-market opportunities and assortative mating causally amplify household inequality and intergenerational persistence? The mathematical statement above fixes the target, assumptions and scope of this catalogue formulation.

\subsection{Sources}
\begin{itemize}\raggedright
\item[\textnormal{[S1]}]\hypertarget{source:30007649:1}{} Benjamin Goldman, Jamie Gracie, Sonya R. Porter. 2026. Who Marries Whom? The Role of Segregation by Race and Class. Abstract, paragraphs on childhood exposure and the spatial marriage-market model.
\url{https://www.nber.org/papers/w35140}.
DOI: 10.3386/w35140.
{\small\textit{Supporting or current literature; see evidence qualification. The paper identifies exposure effects on marriage and quantifies segregation; the full downstream inequality and descendant-mobility problem is not explicitly posed as open.}}
\end{itemize}
\end{document}
